\section{Equality, first variations, and a causal stationary field}
\label{n3:variation:section}

We return to the original exterior in Theorem~\ref{n3:intro:rigidity}.
Write its nonempty compact boundary as
\[
 S=\bigsqcup_{i=1}^{\ell}S_i,\qquad 1\leq\ell<\infty,
\]
where each $S_i$ is a closed connected orientable smooth surface.
No genus is prescribed.  Each component is a future MOTS, and $S$
is outer area-minimizing.  We use the precise outermostness condition
of that theorem: there is no enclosing cut with at least one component
in $\operatorname{Int}\Mext$, every other component either in that
interior or equal to an $S_i$, and $\theta_+\leq0$ on every component.
Thus the condition permits comparison cuts that retain some original
boundary components.  All the arguments in this section
take place on this original exterior.  We use the exterior inequality
of Theorem~\ref{n3:intro:exterior-inequality}, already proved in the
preceding sections, for nearby data on the same manifold with boundary.
Those nearby data need not have an outermost or outer area-minimizing
boundary.  This distinction is essential for the variations below.

Write
\[
 A=|S|_g=16\pi m^2,\qquad
 b^0=\frac Em,\qquad b^i=-\frac{P_i}{m},\qquad
 c=\frac1{2m},\qquad \tau=\tr_g K.
\]
Thus $A=\sum_{i=1}^{\ell}|S_i|_g$ counts all boundary components.
The vector $(b^0,b^i)$ is future unit timelike.  For a nearby ADM
vector define the fixed linear functional
\begin{equation}
 \mathcal E=b^0E_{\rm new}+b^iP_{{\rm new},i}.
 \label{n3:variation:supporting-energy}
\end{equation}
The reverse Lorentzian Cauchy--Schwarz inequality gives
$\mathcal E\ge\sqrt{E_{\rm new}^2-|P_{\rm new}|^2}$ when the
nearby ADM vector is future timelike.  At the original data both sides
equal $m$ and have the same first derivative.  Moreover,
\begin{equation}
 \left.16\pi\frac{\dd}{\dd a}\sqrt{\frac a{16\pi}}
 \right|_{a=A}=c.
 \label{n3:variation:area-coefficient}
\end{equation}

We shall prove the following statement.  The notation $O_2$ for the
vector field uses the original end coordinates, component by component.

\begin{proposition}[Causal adjoint at equality]
\label{n3:variation:adjoint}
There are a function $u$ and a vector field $X$, smooth on the
one-sided manifold $\Mext$, with the following properties.
\begin{enumerate}
 \item $u\ge |X|_g$ everywhere and $u>0$ in $\operatorname{Int}\Mext$.
 For every $q'$ with $1/2<q'<\min\{q,1\}$,
 \begin{equation}
  u=b^0+O_2(r^{-q'}),\qquad
  X=b^i\partial_i+O_2(r^{-q'}).
  \label{n3:variation:asymptote}
 \end{equation}
 \item In the interior,
 \begin{align}
  \sym\nabla X&=-uK,\label{n3:variation:kid-first}\\
  \Delta u&=-\Div(K(X,\cdot)^\sharp),
    \label{n3:variation:lapse-divergence}\\
  \Hess u
   &=u(\Ric_g+\tau K-2K^2)-\mathcal L_XK
       +8\pi X_{(i}J_{j)}.
    \label{n3:variation:kid-second}
 \end{align}
 Here $(K^2)_{ij}=K_i{}^kK_{kj}$, and parentheses denote averaged
 symmetrization.  These equations extend smoothly to $S$.
 \item On $\R\times\operatorname{Int}\Mext$ the Lorentzian metric
 \begin{equation}
  \mathbf g=-u^2\dd t^2
       +g_{ij}(\dd x^i+X^i\dd t)(\dd x^j+X^j\dd t)
  \label{n3:variation:stationary-metric}
 \end{equation}
 has stationary Killing field $\xi=\partial_t$ and induces $(g,K)$
 on $t=0$, with future normal $n=u^{-1}(\partial_t-X)$ and the
 second fundamental form convention $K=\frac12\mathcal L_n g$.
 Its Einstein tensor satisfies
 \begin{equation}
  u\mathbf G_{ij}=-8\pi X_{(i}J_{j)},\qquad
  u\mu+J(X)=0.
  \label{n3:variation:einstein-spatial}
 \end{equation}
 Set $N=u^2-|X|^2$.  Then $\mu N=0$, and, throughout this stationary
 spacetime,
 \begin{equation}
  \mathbf G=
    \frac{8\pi\mu}{u^2}\,\xi^\flat\otimes\xi^\flat,
  \qquad \Scal_{\mathbf g}=0,
  \qquad \Ric_{\mathbf g}(\xi,\cdot)=0.
  \label{n3:variation:null-dust}
 \end{equation}
 In particular, the spacetime is vacuum on $\{N>0\}$.
 \item With $\nu$ pointing into the exterior,
 \begin{equation}
  X=u\nu,\qquad
  \partial_\nu u+K(X,\nu)=\kappa,\qquad
  \kappa=\frac1{4m}
  \quad\hbox{on }S.
  \label{n3:variation:boundary-data}
 \end{equation}
 For each $i=1,\ldots,\ell$, either $u>0$ everywhere on $S_i$,
 or $u=0$ everywhere on $S_i$.  This alternative is made separately
 on each component; the constant $\kappa$ is common to all components.
\end{enumerate}
\end{proposition}

Equation~\eqref{n3:variation:null-dust} deliberately allows nonzero
matter where $N=0$.  This section does not infer vacuum from the
existence of a causal stationary field alone.  The treatment of that
null region belongs to Section~\ref{n3:static:section}.

\subsection{The derivative of the enclosing area}
\label{n3:variation:area-subsection}

We use filled enclosing sets, with coincidence with any component
of $S$ allowed.  To make the perimeter convention explicit for arbitrary
genus, attach a compact orientable handlebody to the inner side of each
$S_i$, identifying its boundary smoothly with $S_i$.  Every closed
connected orientable surface bounds such a handlebody.  The resulting
manifold has no boundary and retains the original asymptotic end.
Extend $g$ smoothly and positively through the collars and over these
caps, by extending its collar coefficients and then joining to a
positive metric on each cap by a partition of unity.  The finite union
of the caps is a compact smooth obstacle with boundary $S$.

Minimize the full perimeter in this capped manifold among relatively
compact finite-perimeter sets containing the obstacle almost everywhere.
No condition on complementary pockets is imposed on these competitors,
so this class is closed under unions and intersections.  A minimizing
set has no bounded complementary pocket: filling one removes its frontier
and decreases perimeter.  Full perimeter counts the frontier coinciding
with $S$, in agreement with Definition~\ref{n3:intro:exterior-class}.
Every contributing boundary lies
in the original closed exterior, so neither the choice of cap nor its
interior metric changes its area.  Only the metric is extended; the
constraints and all variations below remain on $\Mext$.
Let $\mathcal T$ be the collection of boundaries attaining the enclosing
infimum for the original metric.  Each boundary in $\mathcal T$ has
area $A$; $S$ itself is one of them.  Perimeters and tangent-plane
measures count every boundary component.

\begin{lemma}[Compact minimizing cuts and the area envelope]
\label{n3:variation:area-envelope}
Let $g_s$ be a smooth path of metrics through $g$ whose first and
second parameter derivatives are uniformly bounded relative to $g$,
and whose end geometry has uniformly bounded scaled first derivatives
for $s$ near zero.  Put $h=\dot g_0$, and let $a(s)$ be its filled
enclosing infimum.  Then
\begin{equation}
 a'(0+)=\min_{T\in\mathcal T} a'_T(h),\qquad
 a'_T(h)=\frac12\int_T\tr_T h\dd A_g.
 \label{n3:variation:area-envelope-formula}
\end{equation}
The space $\mathcal T$ is compact for the convergence of its sets in
local $L^1$ and of its unoriented tangent-plane area measures.  The
function $T\mapsto a'_T(h)$ is continuous for that topology.
Off the obstacle, two cuts in $\mathcal T$ have the same tangent
plane at each point where they meet.
\end{lemma}

\begin{proof}
A fixed enclosing competitor bounds all minimizing perimeters from
above.  A minimizing boundary that reaches arbitrarily large end
radius has a point with a ball of radius comparable to that radius
which is disjoint from the obstacle.  The scaled metric bounds and the
interior perimeter density estimate give a lower area bound equal to
a positive constant times the square of that radius.  This contradicts
the fixed upper bound.  Thus all minimizing boundaries for the paths
under consideration stay in a fixed compact set; this is also the
compactness argument in Lemma~\ref{n3:reduction:area-compactness}.
One may minimize on a large compact truncation, whose sphere is a
strict mean-curvature barrier, and then remove the truncation.
Perimeter compactness gives a limiting filled set.  The fixed-obstacle
condition is closed under this convergence, and lower
semicontinuity gives a minimizing boundary.

The regularity input is the smooth-obstacle perimeter theorem:
in ambient dimension three, a perimeter minimizer outside a $C^2$
obstacle is $C^{1,1}$ and is smooth minimal away from contact
\cite[Regularity Theorem 1.3(iii)]{HuiskenIlmanen2001}.
These minimizers realize the smooth-cut infimum as well.  Every smooth
filled enclosing cut is an admissible finite-perimeter competitor.
Conversely, flow a regular minimizing boundary for a small positive
time by a smooth vector field pointing strictly outward along $S$ and
supported in its collars.  The moved filled set contains the obstacle
in its interior.  A sufficiently close smooth $C^1$ approximation of
that compact embedded boundary still encloses the obstacle, and its
area tends to the original perimeter as the flow time and approximation
error tend to zero.  Filling any bounded complementary pockets only
decreases area.  The resulting smooth cuts prove equality of the two
infima.  This construction applies to each nearby metric as well.

The local obstacle estimates and the interior density estimates apply uniformly
on the fixed compact region.  For a convergent sequence of minimizing
sets, lower semicontinuity and the upper bound from a fixed minimizing
competitor give convergence of the perimeters.  Express the normal
vector measures in a fixed smooth local $g$-orthonormal frame.  These
are smooth matrix multiples of the distributional derivatives of the
set indicators, so they converge weakly, and their Euclidean total
variations are precisely the convergent $g$-perimeters.
Strict convergence of these vector measures gives convergence of the
integrals of every continuous function of position and unit normal;
in particular it gives convergence for every continuous function of
the unoriented tangent plane.  This is the strict-measure continuity
theorem, applied locally and then with a partition of unity
\cite{AmbrosioFuscoPallara2000}.  The same argument works when the
metrics vary: their uniform convergence converts convergence of the
varying perimeters into strict convergence for the fixed metric.

For any rectifiable cut of bounded area the pointwise area-element
expansion gives, uniformly over those cuts,
\begin{equation}
 |T|_{g_s}=|T|_g+s a'_T(h)+O(s^2)|T|_g.
 \label{n3:variation:uniform-area-taylor}
\end{equation}
Testing a fixed member of $\mathcal T$ gives the upper bound in
Equation~\eqref{n3:variation:area-envelope-formula}.  For the opposite
bound, take minimizers $T_s$ for $g_s$.  Any sequence $s\downarrow0$
has a subsequence converging as above to some $T\in\mathcal T$.
Since $|T_s|_g\ge A$, Equation~\eqref{n3:variation:uniform-area-taylor}
gives
\[
 \frac{a(s)-A}{s}\ge a'_{T_s}(h)+O(s)
       \longrightarrow a'_T(h)
       \ge\min_{T'\in\mathcal T}a'_{T'}(h).
\]
This proves the derivative formula and the asserted continuity.

Finally, if $D_1,D_2$ are two minimizing filled sets, then their union
and intersection are admissible and perimeter submodularity gives
\[
 P(D_1\cup D_2)+P(D_1\cap D_2)
       \le P(D_1)+P(D_2)=2A.
\]
Each term on the left is at least $A$, so both are minimizers.
Different tangent planes at an intersection off the obstacle would
give a corner in one of these two minimizing boundaries, contradicting
their interior regularity.  Thus the tangent planes agree there.
\end{proof}

For later use, the common plane is a continuous field on the union
of these cuts, locally away from $S$.  Indeed, take points $x_j$ on
cuts $T_j$ converging to a point $x$ away from $S$.  Compactness gives
a limiting cut.  The interior density estimate puts $x$ on that cut,
and strict perimeter convergence to its smooth boundary gives local
tangent convergence by the interior regularity theorem.  The common
plane property makes the limit independent of the selected cut.

\subsection{A positive separator for the active constraints}
\label{n3:variation:separation-subsection}

Use the fixed unit ball bundle
\[
 \mathcal B=\{(x,v):x\in\Mext,\ |v|_g\le1\},\qquad
 C(x,v)=16\pi\bigl(\mu(x)+J_x(v)\bigr),
\]
and its closed active subset $\mathcal A=\{C=0\}$.
For a varying metric transport $v$ isometrically by the symmetric
positive square root: if $g_s(\cdot,\cdot)=g(A_s\cdot,\cdot)$,
put $v_s=A_s^{-1/2}v$.  Thus
\begin{equation}
 \dot v_0=-\tfrac12 h^\sharp v,
 \qquad |v_s|_{g_s}=|v|_g.
 \label{n3:variation:frame-derivative}
\end{equation}
The linear variation $C'_H$ below includes this argument variation.

Fix the strict conformal direction from
Lemma~\ref{n3:reduction:strict-direction}, and write it as
\[
 Q=(4\phi g,2\phi K).
\]
Here $\partial_\nu\phi=-1$, $\phi=O_2(r^{-1})$ has a finite flux,
and for some $0<\delta<\min\{q,1\}$ and a smooth positive weight
$w=r^{-3-\delta}$ on the end,
\begin{equation}
 -\Delta\phi-|K||\nabla\phi|\ge w,
 \qquad \frac{-\Delta\phi}{w}\longrightarrow1.
 \label{n3:variation:strict-weight}
\end{equation}
Decrease $\delta$ if necessary in using that lemma.
Let $\mathcal V$ be the real linear space generated by $Q$, all smooth
compactly supported variations $(h,p)$ up to $S$, and the following
four smooth variations cut off to vanish on a fixed large compact
set:
\begin{equation}
 h^{(0)}_{ij}=\frac{\delta_{ij}}r,\qquad p^{(0)}=0;
 \qquad h^{(a)}=0,\qquad p^{(a)}=B(e_a),
 \label{n3:variation:prototypes}
\end{equation}
where, on the end,
\begin{equation}
 B(p)_{ij}=r^{-2}
  \bigl(p_i n_j+p_j n_i-(\delta_{ij}-n_in_j)p\cdot n\bigr),
 \qquad n=x/r.
 \label{n3:variation:momentum-prototype}
\end{equation}
These tensors are Euclidean tracefree and divergence-free.
The scalar-curvature linearization at the Euclidean metric annihilates
$\delta/r$ outside the cutoff.  Consequently the non-strict generators
have
\begin{equation}
 C'_H=O(r^{-3-q})=o(w)
 \quad\hbox{on the end, uniformly for }|v|\le1.
 \label{n3:variation:prototype-error}
\end{equation}
The same estimate includes the frame term in
Equation~\eqref{n3:variation:frame-derivative}.  The momentum prototypes
have independent momentum fluxes and the scalar prototype has nonzero
energy flux.

On an active ray the background conformal scaling term vanishes, so
the exact conformal formulas give
\begin{equation}
 C'_Q=-8\Delta\phi+8K(v,\nabla\phi),\qquad
 \frac{C'_Q}{w}\longrightarrow8,
 \qquad -\theta'_Q=4\quad\hbox{on }S.
 \label{n3:variation:strict-active}
\end{equation}
In particular, every $H\in\mathcal V$ has a well-defined coefficient
$a(H)$ of $Q$, characterized equivalently by
$C'_H/w\to8a(H)$ on the end.  This also shows that the coefficient
is independent of its representation by generators.

\begin{lemma}[Multiplier identity]
\label{n3:variation:separation}
There are a probability measure $\omega$ on $\mathcal T$, a
nonnegative finite measure $\eta$ on $S$, a nonnegative locally finite
measure $\Lambda$ on $\mathcal A$, and a number $z\ge0$ such that,
for every $H\in\mathcal V$,
\begin{equation}
 \begin{split}
 16\pi\mathcal E'(H)
   -c\int_{\mathcal T} a'_T(h)\dd\omega(T)
  =\langle C'_H,\Lambda\rangle
       -\int_S\theta'_H\dd\eta+z a(H).
 \end{split}
 \label{n3:variation:multiplier-identity}
\end{equation}
The pairing is well defined on $\mathcal V$.  The last term vanishes
on compact variations and on the four prototypes.
\end{lemma}

\begin{proof}
Adjoin an end point to $\mathcal A$, collapsing every ray that escapes
spatial compact sets to that point; if $\mathcal A$ is compact, take
an additional isolated end point.  The bounded fibers make the
result a compact space.  Equations~\eqref{n3:variation:prototype-error}
and~\eqref{n3:variation:strict-active} show that $C'_H/w$ extends
continuously to it, with value $8a(H)$.  Take the disjoint compact
union of this space, $S$, and $\mathcal T$, and associate to $H$ the
continuous function whose values on the three pieces are
\begin{equation}
 \frac{C'_H}{w},\qquad -\theta'_H,\qquad
 c a'_T(h)-16\pi\mathcal E'(H),
 \label{n3:variation:test-map}
\end{equation}
respectively.  Smooth dependence of the boundary expansion and
Lemma~\ref{n3:variation:area-envelope} give the remaining continuity.

We first prove that no image of this linear map is strictly positive
at every point of the compact test space.  Suppose otherwise, and
write $H=aQ+H_0$.  Positivity at the end point gives $a>0$.
Realize its tangent by the actual path
\begin{equation}
 g_s=e^{4as\phi}(g+s h_0),\qquad
 K_s=e^{2as\phi}(K+s p_0).
 \label{n3:variation:feasible-path}
\end{equation}
The metric remains positive and uniformly comparable to $g$ for
small $|s|$.  The varied data retain the required asymptotic decay
with exponent $\min\{q,1\}>1/2$, and their charge variations have finite
limits.  The conformal factor has its stated flux.  Thus the path
has differentiable ADM fluxes.

Here are the estimates which ensure nonlinear feasibility.  On the
end the tensors in $H_0$ are $O_2(r^{-1})$ and $O_1(r^{-2})$, and
their far-end Euclidean linear constraints vanish.  Terms in their
linear constraint variation containing a background error cost
$O(s r^{-3-q})$, while new quadratic terms cost $O(s^2r^{-4})$.
Changing the orthonormal identification has the same bound: the
background momentum is $O(r^{-2-q})$, and the metric variation is
$O(r^{-1})$.  Applying the exact conformal law to the intermediate
data $(g+s h_0,K+s p_0)$, and keeping the nonnegative background
term with its positive scaling factor, therefore gives, uniformly
over the full unit ball of test vectors,
\begin{equation}
 e^{4as\phi}C_s
  =C+s\bigl(8a w+o(w)\bigr)+O(s^2w)
  \quad\hbox{on the end}.
 \label{n3:variation:nonlinear-end}
\end{equation}
The $o(w)$ is uniform as $r\to\infty$ at fixed $H$, and the
quadratic bound is uniform for small $s$.  The gradient square in
the conformal formula is $O(r^{-4})=O(w)$ because $\delta<1$.
One may use the composition of the two natural isometries of the
unit balls in deriving this estimate; it tests the same DEC.
At active rays its derivative agrees with
Equation~\eqref{n3:variation:frame-derivative}, since the difference of
two isometric identifications is tangent to the unit sphere and is
annihilated by $J$ at an active ray.

It follows from Equation~\eqref{n3:variation:nonlinear-end} that the DEC
holds on a fixed far end for all sufficiently small positive $s$.
On the remaining compact ball bundle, the derivative is strictly
positive on the closed active set.  A neighborhood of that set has
the same positive derivative with a fixed smaller margin.  On its
compact complement the original $C$ has a positive minimum.
Taylor's formula now gives the DEC throughout the compact region
as well.  The strict inequality $-\theta'_H>0$ on compact $S$ gives
$\theta_{+,s}<0$ for small positive $s$.  The new sources are
integrable: outside a compact set their changes, after the harmless
positive rescaling of the original sources, have the integrable
bounds just displayed.  The ADM vector stays future timelike.
Hence these are data to which
Theorem~\ref{n3:intro:exterior-inequality} applies.

Finally, strict positivity on the compact cut space gives
\[
 c\min_{T\in\mathcal T}a'_T(h)-16\pi\mathcal E'(H)>0.
\]
By Lemma~\ref{n3:variation:area-envelope} and
Equation~\eqref{n3:variation:area-coefficient}, this says that the
right derivative of
$16\pi(\mathcal E(s)-\sqrt{a(s)/(16\pi)})$ is strictly negative.
The expression is zero at $s=0$, whereas the supporting-energy
inequality and Theorem~\ref{n3:intro:exterior-inequality} make it
nonnegative for small positive $s$.  This is a contradiction.

The image of the linear test map is therefore disjoint from the
open cone of strictly positive continuous functions.  The separating
form of the Hahn--Banach theorem gives a nonzero continuous linear
functional, nonnegative on nonnegative functions, which annihilates
the image.  The Riesz representation theorem represents it by
nonnegative finite measures on the three compact pieces.  Its mass
on the objective piece $\mathcal T$ cannot be zero: otherwise its
value on the image of $Q$, which is strictly positive on both
remaining pieces including their end point, would be positive.
Normalize this objective mass to one.  On the finite active rays
divide the corresponding measure by $w$ and call the result
$\Lambda$.  It is locally finite; the original finite measure makes
its pairing with $C'_H$ finite.  Put the end-point mass, multiplied
by eight, into $z$.  Rearranging the annihilation identity gives
Equation~\eqref{n3:variation:multiplier-identity} with the stated signs.
\end{proof}

Define the scalar and vector moments of $\Lambda$ with the fixed
volume-density convention by
\begin{equation}
 \langle u,f\rangle=\int f(x)\dd\Lambda(x,v),\qquad
 \langle X,\alpha\rangle=\int\alpha_x(v)\dd\Lambda(x,v).
 \label{n3:variation:moments}
\end{equation}
Initially these are measures, including possible measures on $S$.
Their positivity and the support of $\Lambda$ give, distributionally,
\begin{equation}
 u\ge |X|_g,\qquad u\mu+J(X)=0.
 \label{n3:variation:measure-complementarity}
\end{equation}
The first statement means domination of the total variation of the
vector measure by the scalar measure; in particular it can be tested
against arbitrary continuous unit covector fields.

\subsection{Removing the interior area multipliers}
\label{n3:variation:normal-slice-subsection}

\begin{lemma}[Normal slice tests]
\label{n3:variation:normal-tests}
For every smooth $s$ compactly supported in
$\operatorname{Int}\Mext$ there are compact variations $H_\epsilon$
with metric component $h_\epsilon=2sK$ such that
\begin{equation}
 C'_{H_\epsilon}\longrightarrow0
 \quad\hbox{uniformly on the active rays over compact sets}.
 \label{n3:variation:normal-test-limit}
\end{equation}
Their supports lie in one fixed compact set.
\end{lemma}

\begin{proof}
Choose a compact neighborhood of the support of $s$.  In a Gaussian
time collar put
\[
 \mathbf g_\epsilon=-\dd t^2+g_\epsilon(t),\qquad
 g_\epsilon(0)=g,\qquad \partial_tg_\epsilon(0)=2K.
\]
Choose the second time derivative so that its spatial Einstein
components at time zero are
\begin{equation}
 (\mathbf G_\epsilon)_{ij}=8\pi(D_\epsilon)_{ij},
 \qquad D_\epsilon=\frac{J\otimes J}{\mu+\epsilon}.
 \label{n3:variation:dust-jets}
\end{equation}
This is a free choice of jets, not a choice of evolution satisfying
an energy condition.  To see its solvability, the terms in the
spatial Einstein tensor containing $\partial_tK$ are its trace
reversal, with a fixed nonzero overall sign.  On symmetric tensors
in dimension three this map is invertible: its eigenvalue is one
on the tracefree part and minus two on the pure-trace part.
All remaining terms are determined by $g,K$ and their spatial
derivatives.  Thus smooth second jets realizing
Equation~\eqref{n3:variation:dust-jets} exist; a quadratic extension
with a time cutoff realizes them by a smooth Lorentzian metric for
sufficiently small time.  The constraints give at time zero
\begin{equation}
 \mathbf G_\epsilon(n,n)=8\pi\mu,
 \qquad \mathbf G_\epsilon(n,e_i)=8\pi J_i,
 \label{n3:variation:gaussian-constraints}
\end{equation}
where $n=\partial_t$ and $e_i$ is a spatial orthonormal frame.

The regularization has the precise compact convergence needed below.
Define $D=J\otimes J/\mu$ where $\mu>0$, and $D=0$ at $\mu=0$.
The DEC implies
\begin{equation}
 |D_\epsilon|\le\mu,\qquad
 |\nabla D_\epsilon|\le2|\nabla J|+|\nabla\mu|.
 \label{n3:variation:dust-c1-bound}
\end{equation}
The second inequality follows by differentiating the quotient and
using $|J|\le\mu$ in each term.  Smooth nonnegative $\mu$ has
$\nabla\mu=0$ at its zero set.  There also $\nabla J=0$:
in a smooth frame $|J_i|\le\mu$, so each directional derivative of
$J_i$ vanishes.  It follows first that $D$ is differentiable with
zero derivative at that set, and then from
Equation~\eqref{n3:variation:dust-c1-bound} that this derivative is
continuous there.  On a compact set, split into $\mu\ge\delta$ and
$\mu<\delta$.  On the first part the quotient and its derivatives
converge uniformly as $\epsilon\to0$.  On the second part the bound
in Equation~\eqref{n3:variation:dust-c1-bound} is uniformly small as
$\delta\downarrow0$, by compactness and continuity at the zero set.
Therefore
\begin{equation}
 D_\epsilon\longrightarrow D\quad\hbox{in }C^1
 \hbox{ on every compact set}.
 \label{n3:variation:dust-c1-convergence}
\end{equation}

Vary the slice by the graphs $t=a s(x)$ at $a=0$, using its induced
metric and second fundamental form.  Their variation is compactly
supported where $s$ and its derivatives are supported, and its metric
component is $2sK$.  We prove that its first DEC variation tends to
zero at an active ray $(x,v)$.

If $\mu(x)>0$, activity implies
$\mu=|J|$, $|v|=1$, and $v=-J^\sharp/\mu$ at that point.
Write $\zeta=n+v$.  On the portion of the initial slice where
$\mu>0$ the limiting tensor prescribed above is
\begin{equation}
 \mathbf G_0=\frac{8\pi}{\mu}\alpha\otimes\alpha,
 \qquad \alpha(n)=\mu,\qquad \alpha(e_i)=J_i.
 \label{n3:variation:causal-rank-one}
\end{equation}
The covector $\alpha$ is causal and is nonnegative on future causal
vectors.  At the active point, $\zeta$ is null and
$\alpha(\zeta)=0$.  Parallel transport $\zeta$ along any spatial
curve in the initial slice, using the spacetime connection.  It
remains future null, so its contraction with $\alpha$ is a
nonnegative function with a minimum zero at $x$.  Thus
\[
 (\nabla_i\alpha)(\zeta)=0,\qquad
 (\nabla_i\mathbf G_0)(e_i,\zeta)=0
 \quad\hbox{at }x.
\]
The spatial spacetime connection at time zero depends only on $g,K$,
not on $\epsilon$.  Equation~\eqref{n3:variation:dust-c1-convergence}
therefore also gives convergence of these spatial covariant
derivatives.  Contracted Bianchi, applied before taking the limit,
states in an orthonormal frame that
\begin{equation}
 (\nabla_n\mathbf G_\epsilon)(n,\zeta)
   =\sum_i(\nabla_{e_i}\mathbf G_\epsilon)(e_i,\zeta).
 \label{n3:variation:bianchi-normal}
\end{equation}
Only the normal derivative of the constraint components is used here;
Bianchi expresses it through the spatial first derivatives controlled
above.
Its limiting value is zero.

The derivative of the contraction defining the DEC on the graph
also differentiates its two vector arguments.  The term varying
the first argument vanishes because
$\mathbf G_0(\cdot,\zeta)=0$.  For the second argument, its
isometrically transported spatial part still has length one, so the
varied $\zeta$ stays null.  At the active point $\alpha$ is
proportional to the covector obtained by lowering $\zeta$; hence
$\alpha(\dot\zeta)=0$.  This kills the second argument term.
Together with Equation~\eqref{n3:variation:bianchi-normal} it proves
the desired limit at such a ray.

If $\mu(x)=0$, then $J=0$ and their spatial first derivatives vanish.
The limiting tensor and its spatial first derivatives are zero.
Equation~\eqref{n3:variation:bianchi-normal} now kills the required
normal derivative, and both argument terms vanish.  This applies to
all active $v$, including $|v|<1$.  Finally, all the expressions just
used involve only fixed compactly bounded jets, $D_\epsilon$, and
its spatial first derivatives.  The compact $C^1$ convergence makes
the convergence uniform, including near $\mu=0$.
This proves Equation~\eqref{n3:variation:normal-test-limit}.
\end{proof}

\begin{lemma}[Concentration of the area multiplier]
\label{n3:variation:area-support}
The measure $\omega$ in Lemma~\ref{n3:variation:separation} is
concentrated on the single cut $S$.
\end{lemma}

\begin{proof}
Use the variations of Lemma~\ref{n3:variation:normal-tests} in
Equation~\eqref{n3:variation:multiplier-identity}.  Their ADM and
boundary variations vanish.  Their supports lie in a fixed compact
set, on which $\Lambda$ is finite, so uniform convergence permits
passage to the limit in its pairing.  The metric variation is $2sK$,
independent of the regularization.  It follows that
\begin{equation}
 \int_{\mathcal T}\int_T s\tr_TK\dd A_g\dd\omega(T)=0
 \quad\hbox{for every }s\in C_c^\infty(\operatorname{Int}\Mext).
 \label{n3:variation:aggregate-trace}
\end{equation}

Consider the positive aggregate area measure
$\beta=\int_{\mathcal T}\dd A_T\dd\omega(T)$ off $S$.
By Lemma~\ref{n3:variation:area-envelope} and the observation following
its proof, the function $F(x)=\tr_{\Pi(x)}K(x)$, where $\Pi(x)$ is
the common tangent plane, is continuous on the union of the cuts,
locally away from $S$.  In particular it is measurable there.
Equation~\eqref{n3:variation:aggregate-trace} says the signed Radon
measure $F\beta$ is zero.  Thus $F=0$ $\beta$-almost everywhere.
Fubini's theorem on a countable compact exhaustion shows that for
$\omega$-almost every cut its tangential trace of $K$ is zero at
area-almost every point off $S$.  Each such cut is smooth minimal
there, so continuity improves this to
\[
 H_T=0,\qquad \tr_TK=0\quad\hbox{on }T\setminus S.
\]
There can be no cancellation in this argument between cuts with
different tangent planes: the common-plane property is exactly what
made $F$ a single function of position.

Fix one of these cuts.  Near a point of contact with $S$, the
$C^{1,1}$ obstacle graph and the smooth graph of $S$ have the same
first derivatives on their coincidence set.  Their second
derivatives agree almost everywhere on that set.  One way to see
the latter assertion is to apply the fact that a Lipschitz function
vanishing on a measurable set has derivative zero almost everywhere
on that set, first to the graph difference and then to each of its
first derivatives.  The outward orientations agree at contact.
Since $S$ is a MOTS, the cut consequently satisfies
$H_T+\tr_TK=0$ almost everywhere across contact, as well as on its
complement.  In a graph chart this is a uniformly elliptic
quasilinear equation for a $C^{1,1}$ function.  Its coefficients are
Hölder in position and first derivatives.  Interior elliptic
regularity followed by differentiation therefore makes the graph
smooth.  Thus the whole cut is a smooth MOTS.

Suppose the cut touches a component $S_i$.  At each contact point,
the two smooth graph equations have the same orientation and one graph
lies on one side of the other.  Their difference satisfies a linear
uniformly elliptic equation with bounded coefficients.  The strong
comparison principle implies local coincidence.  The contact set
$T\cap S_i$ is therefore open in $S_i$; it is also closed by compactness.
Connectedness of $S_i$ gives $S_i\subset T$.  Local coincidence and
compactness then make $S_i$ a connected component of $T$.  Consequently,
every component of $T$ either equals one of the original $S_i$ or is
disjoint from $S$ and lies entirely in $\operatorname{Int}\Mext$.
This argument does not use the genus of $S_i$.

The cut is smooth and encloses the entire obstacle.  Its filled set
has no bounded complementary pocket, since filling such a pocket would
remove positive perimeter and contradict minimality.  If $T$ had any
interior component, it would thus be an enclosing smooth MOTS with at
least one interior component and with every remaining component either
interior or equal to an original $S_i$.  This is exactly a cut excluded
by the partial-coincidence outermostness hypothesis, even if it retains
some original boundary components.  Hence $T$ has no interior component.
There is a subset $I\subset\{1,\ldots,\ell\}$ such that
\[
 T=\bigsqcup_{i\in I}S_i,\qquad
 \sum_{i\in I}|S_i|_g=|T|_g=A
                       =\sum_{i=1}^{\ell}|S_i|_g.
\]
Each $|S_i|_g$ is strictly positive.  The equality of these finite sums
therefore forces $I=\{1,\ldots,\ell\}$, so $T=S$.
It follows that $\omega$-almost every cut is $S$, proving the lemma.
\end{proof}

\subsection{The interior adjoint equations}
\label{n3:variation:interior-subsection}

We now know that the area term in
Equation~\eqref{n3:variation:multiplier-identity} is exactly $c a'_S(h)$.
In particular, arbitrary variations compactly supported in the open
exterior have zero total multiplier pairing.  We first use this fact
before making any regularity assumption on the moments.

\begin{lemma}[Interior regularity and the full metric equation]
\label{n3:variation:interior-equations}
The moments $(u,X)$ are smooth in the open exterior and satisfy
Equations~\eqref{n3:variation:kid-first}--\eqref{n3:variation:kid-second}
and $u\mu+J(X)=0$ there.  Their full first jet satisfies a homogeneous
linear first-order system with smooth background coefficients.
\end{lemma}

\begin{proof}
For a compact variation $p=\dot K$ with $h=0$, direct differentiation
and integration of the momentum divergence give the coefficient
\[
 2\bigl[u(\tau g-K)-\sym\nabla X+(\Div X)g\bigr]
\]
against $p$.  It is zero distributionally.  Taking its trace gives
$\Div X=-u\tau$, and substituting back proves
Equation~\eqref{n3:variation:kid-first}.

Next use the compact simultaneous conformal variation
$(h,p)=(4\psi g,2\psi K)$.  Its background scaling term is
$-64\pi\psi(u\mu+J(X))$, which vanishes by
Equation~\eqref{n3:variation:measure-complementarity}.  The remaining
pairing is
\[
 -8\langle u,\Delta\psi\rangle
     +8\langle X,K(\nabla\psi,\cdot)\rangle=0.
\]
This is Equation~\eqref{n3:variation:lapse-divergence} in distributions.
Taking a divergence of the symmetric-gradient equation and using
$\Div X=-u\tau$ also gives
\begin{equation}
 \begin{split}
  \Delta X_j+\Ric_{jk}X^k
    &=(\tau\delta_j{}^i-2K_j{}^i)\nabla_i u
           +u(\nabla_j\tau-2\nabla^iK_{ij}).
 \end{split}
 \label{n3:variation:elliptic-shift}
\end{equation}
Together with the lapse equation this is a second-order elliptic
system with scalar Laplacian principal part and smooth lower-order
couplings.  Locally the measure moments belong to some negative
Sobolev space.  The local Laplace estimate applied to this system
improves their Sobolev order by one, since the right sides have at
most one derivative of the unknowns.  Iterating with nested compact
cutoffs puts them in every local Sobolev space.  Sobolev embedding
then gives smoothness.  The distributional inequalities and
complementarity now hold pointwise in the interior.

We give the full metric calculation, including the dependence of
the momentum test vectors on the metric.  Put $P=K-\tau g$.
For compactly supported variations the constraint pairing is the
first variation of
\begin{equation}
 I(g,K)=\int
  \left[u(R+\tau^2-|K|^2)+2X^i\nabla^jP_{ij}\right]\dd V_g,
 \label{n3:variation:constraint-action}
\end{equation}
together with the frame correction
\begin{equation}
 -8\pi\int h(X,J^\sharp)\dd V_g.
 \label{n3:variation:frame-correction}
\end{equation}
The integration may be localized to a fixed region containing the
variation.  Changing its volume form instead of retaining the
original fixed measure adds one half of its background integrand
times $\tr h$.  That integrand is $16\pi(u\mu+J(X))=0$,
so this replacement does not change the pairing.
Equation~\eqref{n3:variation:frame-correction} follows directly from
Equation~\eqref{n3:variation:frame-derivative} and has the displayed
negative sign.

Integrating the shift term once by parts makes it
$-\int P^{ij}(\mathcal L_Xg)_{ij}\dd V_g$ up to a boundary term
unaffected by these variations.  In taking a pure metric variation,
$u$, contravariant $X$, and covariant $K$ are fixed.  The elementary
variation formulas used are
\begin{align*}
 \dot R&=-\Ric^{ij}h_{ij}
             +\nabla^i\nabla^jh_{ij}-\Delta(\tr h),\\
 \dot\tau&=-K^{ij}h_{ij},\qquad
 \bigl(|K|^2\bigr)^{\boldsymbol\cdot}
       =-2(K^2)^{ij}h_{ij},\qquad
 (\dd V)^{\boldsymbol\cdot}=\tfrac12\tr h\dd V,\\
 (P^{ij})^{\boldsymbol\cdot}
       &=-h^i{}_aK^{aj}-h^j{}_aK^{ia}
                +(K^{ab}h_{ab})g^{ij}+\tau h^{ij},\\
 (\mathcal L_Xg)^{\boldsymbol\cdot}&=\mathcal L_Xh.
\end{align*}
For example, integrating the last term by parts and combining it
with the variation of $P^{ij}$ gives the following coefficient from
the entire shift integral:
\[
 (\mathcal L_XK)_{ij}-(\Div X)K_{ij}
       -\bigl[X(\tau)+K^{ab}\nabla_aX_b\bigr]g_{ij}.
\]
The scalar-curvature terms contribute
$\Hess u-(\Delta u)g-u\Ric$.  Consequently the vanishing metric
coefficient is the explicit equation
\begin{equation}
 \begin{split}
 0={}&\nabla_i\nabla_j u-(\Delta u)g_{ij}-u\Ric_{ij}
       +2u(K^2-\tau K)_{ij}\\
   &+\frac u2(R+\tau^2-|K|^2)g_{ij}
       +(\mathcal L_XK)_{ij}-(\Div X)K_{ij}\\
   &-\bigl[X(\tau)+K^{ab}\nabla_aX_b\bigr]g_{ij}
       -8\pi X_{(i}J_{j)}.
 \end{split}
 \label{n3:variation:full-metric-adjoint}
\end{equation}
This calculation can equally be read distributionally, since each
coefficient is linear in the moments and their derivatives.

For completeness, its simplification uses no vacuum assumption.
Equation~\eqref{n3:variation:kid-first} gives
\[
 \Div X=-u\tau,\quad
 K^{ab}\nabla_aX_b=-u|K|^2,\quad
 \tr(\mathcal L_XK)=X(\tau)-2u|K|^2.
\]
Taking the trace of Equation~\eqref{n3:variation:full-metric-adjoint}
and using
$16\pi\mu=R+\tau^2-|K|^2$ and $J(X)=-u\mu$ yields
\begin{equation}
 \Delta u+X(\tau)=\frac u2(R+\tau^2+|K|^2).
 \label{n3:variation:trace-metric-equation}
\end{equation}
Substitution of this identity into
Equation~\eqref{n3:variation:full-metric-adjoint} gives exactly
Equation~\eqref{n3:variation:kid-second}.

These equations have the asserted finite-type property.  The
right side of Equation~\eqref{n3:variation:kid-second} is linear in
$u,X,\nabla X$, with smooth background coefficients, since
\[
 (\mathcal L_XK)_{ij}
   =X^k\nabla_kK_{ij}
            +K_{kj}\nabla_iX^k+K_{ik}\nabla_jX^k.
\]
The second derivatives of $X$ are also fixed by first jets.  To
make this explicit, put $B_{ij}=\nabla_{(i}X_{j)}=-uK_{ij}$.
Commuting covariant derivatives gives
\begin{equation}
 \begin{split}
 \nabla_i\nabla_jX_k
   ={}&\nabla_i B_{jk}+\nabla_j B_{ik}-\nabla_k B_{ij}\\
     &+\tfrac12\left(
        [\nabla_i,\nabla_j]X_k
        -[\nabla_i,\nabla_k]X_j
        -[\nabla_j,\nabla_k]X_i\right).
 \end{split}
 \label{n3:variation:prolongation}
\end{equation}
The commutators are curvature times $X$, and differentiating $B=-uK$
uses only $u,\nabla u$ and the background first derivative of $K$.
Thus every derivative of $(u,X,\nabla u,\nabla X)$ is a homogeneous
linear expression in that same first jet.  This proves the final
assertion.
\end{proof}

\subsection{Boundary continuation and the ADM normalization}
\label{n3:variation:continuation-subsection}

\begin{lemma}[Continuation, boundary atoms, and the end constants]
\label{n3:variation:regularity-normalization}
The smooth interior fields extend smoothly to the one-sided boundary.
The original moment measures have no boundary-supported part.  They
satisfy Equation~\eqref{n3:variation:asymptote}, and $u>0$ in the
interior.
\end{lemma}

\begin{proof}
In a smooth normal collar $(s,y)$ of $S$, all background coefficients
in the first-jet system of Lemma~\ref{n3:variation:interior-equations}
are smooth up to $s=0$.  Along a normal line the system is an ordinary
linear differential equation for the full first jet, with bounded
coefficients.  Solve it down to $s=0$ from a fixed interior collar
surface.  Smooth dependence on the initial data and on $y$ supplies
a smooth extension, which agrees with the original solution by
uniqueness along each normal line.  In particular all its one-sided
derivatives have the values supplied by the system.  The same system
also proves that an interior first jet which vanishes at one point
vanishes everywhere on the connected interior, by continuation
along piecewise smooth paths.

We next establish the asymptotic form without having prescribed
growth to the measures.  Let $Y$ denote all coordinate components
of $(u,X)$.  The original $O_2(r^{-q})$, $O_1(r^{-1-q})$ bounds,
the constraints, and Equations~\eqref{n3:variation:kid-second}
and~\eqref{n3:variation:prolongation} give
\begin{equation}
 |\partial^2Y|
    \le C r^{-1-q'}|\partial Y|
          +C r^{-2-q'}|Y|,
 \qquad \tfrac12<q'<\min\{q,1\}.
 \label{n3:variation:end-jet-estimate}
\end{equation}
Only the original derivative bounds occur here: the curvature and
$\partial K,J$ are $O(r^{-2-q'})$, and the coefficients of first
jets are $O(r^{-1-q'})$.

Fix a large sphere.  Radial integration of
Equation~\eqref{n3:variation:end-jet-estimate}, with the supremum of
$|Y|/r+|\partial Y|$ up to radius $r$, bounds this supremum by its
initial value plus its integral against $C r^{-1-q'}\dd r$.
Gronwall's inequality therefore gives, uniformly in angle,
\[
 |Y|=O(r),\qquad |\partial Y|=O(1),\qquad
 |\partial^2Y|=O(r^{-1-q'}).
\]
Each coordinate gradient has a limit along a radial ray, with error
$O(r^{-q'})$.  Two points on a sphere can be joined on that sphere
by a path of length $O(r)$; the last Hessian estimate shows that
their gradient limits agree.  Thus there is one constant matrix $L$
such that
\[
 \partial Y=L+O(r^{-q'}),\qquad
 Y=Lx+O(r^{1-q'}).
\]
The scalar field $u$ is nonnegative on every large sphere, so its
linear coefficient is zero.  Then $u=O(r^{1-q'})$, and
$|X|\le u$ eliminates every linear coefficient of $X$ as well.
Now Equation~\eqref{n3:variation:end-jet-estimate} improves to
$|\partial^2Y|=O(r^{-1-2q'})$.  Since the gradients have limit zero,
radial integration gives $|\partial Y|=O(r^{-2q'})$.  As $2q'>1$,
the values have limits along radial rays.  Their differences on a
large sphere are $O(r^{1-2q'})$, so there is a single constant limit
$Y_\infty$.  With $Y$ bounded, one more use of
Equation~\eqref{n3:variation:end-jet-estimate} gives
\begin{equation}
 Y=Y_\infty+O(r^{-q'}),\qquad
 \partial Y=O(r^{-1-q'}),\qquad
 \partial^2Y=O(r^{-2-q'}).
 \label{n3:variation:unnormalized-asymptote}
\end{equation}

At this point the original measure moments could still have a
component supported on $S$.  Subtract the integration by parts of
the smooth interior fields from the compact identity
\eqref{n3:variation:multiplier-identity}.  Test with $p=0$ and a metric
variation whose value and full first jet vanish on $S$.  The area
and expansion variations and all the smooth integrated boundary
terms then vanish.  The remaining scalar-curvature variation on
$S$ contains
\[
 -\partial_\nu^2(\tr_S h).
\]
This is an arbitrary smooth function on $S$: in collar coordinates
one may take the tangential trace to be a prescribed multiple of
$s^2$ times a cutoff.  All momentum and frame terms involve only
the value or first jet of $h$ and hence vanish on $S$.  It follows
that the scalar moment has no boundary measure.  Domination
$|X|\le u$ as measures eliminates a vector boundary measure also.

We can therefore integrate the smooth adjoint by parts all the way
to $S$ and to a large coordinate sphere.  Use a prototype from
Equation~\eqref{n3:variation:prototypes}, with its cutoff chosen to
vanish near $S$.  Its area and expansion variations vanish, and
$a(H)=0$.  Its constraint pairing is integrable by
Equation~\eqref{n3:variation:prototype-error} and
Equation~\eqref{n3:variation:unnormalized-asymptote}.  Interior terms
vanish by the adjoint equations.  If the constants in the latter
equation are $(u_\infty,X_\infty)$, its outer boundary term is
\begin{equation}
 \begin{split}
 \int_{S_R}\bigl[&u_\infty(\partial_jh_{ij}-\partial_i h_{jj})
       +2X_\infty^j(p_{ij}-(\tr_\delta p)\delta_{ij})\bigr]
          n^i\dd A_\delta+o(1).
 \end{split}
 \label{n3:variation:adjoint-charge-flux}
\end{equation}
For clarity, every omitted term vanishes under the stated decay:
the derivatives of $u,X$ are $O(r^{-1-q'})$ and multiply
$h=O(r^{-1})$; their nonconstant values multiply
$\partial h,p=O(r^{-2})$; and the metric-variation momentum terms
contain $Kh=O(r^{-2-q'})$.  Multiplication by sphere area still
leaves $O(r^{-q'})$.  Thus the limit of
Equation~\eqref{n3:variation:adjoint-charge-flux} is
$16\pi(u_\infty E'+X_\infty^iP_i')$.

The scalar prototype has $E'=1/2$ and $P'=0$; for the momentum
prototype in direction $p$, $E'=0$ and $P'=2p/3$.  The latter follows
by integrating
$B(p)_{ij}n^j=r^{-2}(p_i+(p\cdot n)n_i)$.
Equation~\eqref{n3:variation:multiplier-identity} for these four
independent fluxes consequently gives
\[
 u_\infty=b^0,\qquad X_\infty^i=b^i.
\]
This proves Equation~\eqref{n3:variation:asymptote} and in particular
shows that the adjoint is nontrivial.

If $u$ vanished at an interior point, smooth nonnegativity would
give $\nabla u=0$ there.  The inequality $|X|\le u$ gives $X=0$;
since $u$ has zero differential, it also forces $\nabla X=0$ at
that point.  The whole first jet would be zero.  Its continuation
property would make $(u,X)$ identically zero, contradicting
$u_\infty=b^0>0$.  Thus $u>0$ in the interior.
\end{proof}

\subsection{The stationary Einstein tensor}
\label{n3:variation:stationary-subsection}

We can now form the nondegenerate Lorentzian metric
\eqref{n3:variation:stationary-metric} on the open exterior.  All its
coefficients are independent of $t$.  Its future unit normal to a
constant-time slice is $n=u^{-1}(\partial_t-X)$, and
\[
 \frac12\mathcal L_n g
       =-\frac1{2u}\mathcal L_Xg=K
\]
by Equation~\eqref{n3:variation:kid-first}.  Thus this construction
recovers the original second fundamental form, including its sign.

Here is an explicit curvature verification of the stationary
equation; in particular the modified metric adjoint is not being
identified with vacuum Killing initial data by assertion.  The
Gauss--Codazzi formulas and the stationary normal-derivative formula
for $K$, with the convention just specified, give
\begin{align}
 \Ric_{\mathbf g,ij}
   &=\Ric_{g,ij}+\tau K_{ij}-2(K^2)_{ij}
       -u^{-1}\bigl[(\mathcal L_XK)_{ij}
                             +(\Hess u)_{ij}\bigr],
       \label{n3:variation:stationary-ricci}\\
 \Scal_{\mathbf g}
   &=R+\tau^2+|K|^2
             -2u^{-1}\bigl(\Delta u+X(\tau)\bigr).
       \label{n3:variation:stationary-scalar}
\end{align}
These identities can also be checked from the lapse-shift metric
without an evolution assumption: its scalar decomposition is
\[
 u\Scal_{\mathbf g}
  =u(R+|K|^2-\tau^2)
            -2\Div(\nabla u+\tau X),
\]
and $\Div X=-u\tau$ changes this into
Equation~\eqref{n3:variation:stationary-scalar}.  The spatial Ricci
formula is its tensor counterpart, with normal change
$u^{-1}(-\mathcal L_XK)$ and lapse acceleration $u^{-1}\Hess u$.
Equivalently, variation of the reduced scalar action
$\int u(R+|K|^2-\tau^2)\dd V$ at fixed lapse and contravariant
shift has coefficient $-u\mathbf G_{ij}$.
This agrees with the explicit coefficient
\eqref{n3:variation:full-metric-adjoint} before its frame correction.

Equation~\eqref{n3:variation:trace-metric-equation} in
Equation~\eqref{n3:variation:stationary-scalar} gives
$\Scal_{\mathbf g}=0$.  Equation~\eqref{n3:variation:kid-second} in
Equation~\eqref{n3:variation:stationary-ricci} now gives
\[
 u\mathbf G_{ij}=u\Ric_{\mathbf g,ij}
                  =-8\pi X_{(i}J_{j)}.
\]
The normal and mixed Einstein components are the constraints,
$\mathbf G(n,n)=8\pi\mu$ and
$\mathbf G(n,e_i)=8\pi J_i$.  This proves
Equation~\eqref{n3:variation:einstein-spatial}.

It remains to draw exactly the correct consequence of causality.
At every interior point,
\begin{equation}
 0=u\mu+J(X)
    \ge u\mu-|J||X|
    \ge u(\mu-|J|)\ge0.
 \label{n3:variation:causal-equalities}
\end{equation}
If $u>|X|$, these equalities force $\mu=|J|=0$.
If $\mu>0$, they force
\begin{equation}
 |X|=u,\qquad |J|=\mu,\qquad
 J=-\frac\mu u X^\flat.
 \label{n3:variation:dust-alignment}
\end{equation}
When $\mu=0$, the DEC gives $J=0$ and all Einstein components
already vanish.  When $\mu>0$, use the orthonormal coframe with
$\theta^0(n)=1$ and $\theta^i(e_j)=\delta^i_j$.  The constraints
and the spatial Einstein equation give
\[
 \mathbf G
  =8\pi\mu\left(\theta^0-\frac{X_i}{u}\theta^i\right)^2
  =\frac{8\pi\mu}{u^2}\xi^\flat\otimes\xi^\flat.
\]
The covector in parentheses is null and annihilates
$\xi=un+X$.  This proves the tensor identity in
Equation~\eqref{n3:variation:null-dust} everywhere, with
$\mu N=0$.  Its trace is zero and its contraction with $\xi$ is
zero, so $\Ric_{\mathbf g}(\xi,\cdot)=0$.  All these identities
hold on every stationary time translate.  No strict timelikeness
has been inferred on the internal null set.

\subsection{Free boundary variations}
\label{n3:variation:boundary-subsection}

We finish the proof of Proposition~\ref{n3:variation:adjoint} by
determining the boundary data.  By
Lemma~\ref{n3:variation:regularity-normalization}, all moments now
have smooth volume densities up to $S$, with no constraint measure
on $S$ itself.  The integration-boundary normal is $-\nu$.

Take an arbitrary compact $K$ variation $p$, allowing arbitrary
boundary values, with $h=0$.  Integration of the momentum term has
boundary contribution
\[
 -2\int_S\bigl[p(X,\nu)-(\tr p)X_\nu\bigr]\dd A.
\]
The boundary expansion variation is $\tr_Sp$, and the area variation
vanishes.  Separating the mixed and tangential-trace components of
$p$ in Equation~\eqref{n3:variation:multiplier-identity} gives
\begin{equation}
 X_T=0,\qquad \dd\eta=2X_\nu\dd A.
 \label{n3:variation:eta-density}
\end{equation}
Indeed, the remaining boundary expression is
\[
 -2\int_S p(X_T,\nu)\dd A
      +\int_S(\tr_Sp)(2X_\nu\dd A-\dd\eta),
\]
and these components of $p$ are independent.  This also identifies
$\eta$ as a smooth density.

Now use the compact conformal variation $(4\psi g,2\psi K)$,
allowing arbitrary value and normal derivative of $\psi$ on $S$.
Its area derivative is $4\int_S\psi\dd A$, and its expansion
derivative is $4\partial_\nu\psi$, since $\theta_+=0$.
Integrating the lapse equation by parts with normal $-\nu$ yields
\begin{equation}
 \begin{split}
 -4c\int_S\psi\dd A
    ={}&8\int_S\bigl[u\partial_\nu\psi
       -(\partial_\nu u+K(X,\nu))\psi\bigr]\dd A\\
      &-4\int_S\partial_\nu\psi\dd\eta.
 \end{split}
 \label{n3:variation:conformal-boundary}
\end{equation}
The independently prescribed normal derivative gives
$\dd\eta=2u\dd A$.  Comparing with
Equation~\eqref{n3:variation:eta-density} yields $X=u\nu$.
The independently prescribed value then gives
\[
 \partial_\nu u+K(X,\nu)=\frac c2=\frac1{4m}.
\]
This proves Equation~\eqref{n3:variation:boundary-data} with its
normal orientation and numerical normalization.  The boundary values
and normal derivatives used in these tests can be supported near any
one $S_i$.  The same total-area coefficient $c$ therefore gives the
same constant $\kappa=c/2$ on every boundary component.

Let $L$ be the MOTS stability operator for outward normal graph
variation at $S$: a graph with initial speed $s\nu$ has expansion
derivative $Ls$.  Its principal part is $-\Delta_S$ and its
coefficients are smooth.  Extend an arbitrary smooth $s\nu$ to
a compact vector field $Y$ on the one-sided exterior and test with
\[
 h=\mathcal L_Yg,\qquad p=\mathcal L_YK.
\]
Only its one-sided jets at $S$ are needed; an actual diffeomorphism
of the manifold with boundary is not required to be a member of
the variation space.  In the open exterior these are spatial
diffeomorphism variations.  Their DEC derivative vanishes at every
active ray: under diffeomorphism transport, the derivative is a
spatial derivative of a nonnegative scalar at its zero; the
difference from the isometric vector transport is tangent to the
unit sphere, and $J$ annihilates that difference at an active ray.
At $\mu=0$ the latter statement also holds since $J=0$.
There is no boundary constraint measure left to test.

The ADM variation vanishes, the area variation is
$\int_SHs\dd A$, and the expansion variation is $Ls$.
Using $\dd\eta=2u\dd A$ in
Equation~\eqref{n3:variation:multiplier-identity} gives
\[
 c\int_S Hs\dd A=2\int_S uLs\dd A
 \quad\hbox{for every }s\in C^\infty(S).
\]
Consequently
\begin{equation}
 2L^*u=cH\quad\hbox{on }S.
 \label{n3:variation:boundary-stability}
\end{equation}
Outer area minimization implies $H\ge0$: the first area variation
of every nonnegative outward speed is nonnegative.  Thus the
nonnegative smooth $u|_S$ satisfies $L^*u\ge0$.  The strong minimum
principle applies with no sign assumption on the original
zeroth-order coefficient: subtract a sufficiently large constant
from the operator with positive Laplacian principal part, using
$u\ge0$, to obtain the usual nonpositive zeroth-order coefficient.
Apply this principle on each connected $S_i$.  It follows that
either $u>0$ everywhere on $S_i$, or $u$ is identically zero on $S_i$
\cite[Chapter 3]{GilbargTrudinger2001}.  The argument is componentwise
and does not require the same alternative on distinct components.

All assertions of Proposition~\ref{n3:variation:adjoint} have now been
proved.  On any component $S_i$ where $u=0$, one also has
$X=0$ and $\nabla X=0$: tangential derivatives of $X=0$ vanish,
and Equation~\eqref{n3:variation:kid-first} supplies the remaining
normal derivatives there.  On any component $S_i$ where $u>0$,
the same equation and Equation~\eqref{n3:variation:boundary-data} give
\[
 \partial_\nu N
   =2u\bigl(\partial_\nu u+K(X,\nu)\bigr)=2u\kappa
   \quad\hbox{on }S_i.
\]
These two local boundary behaviors apply independently at each component
and will be used in the static classification.
